\documentclass{article}
\usepackage[margin=1.15in]{geometry}
\usepackage{xcolor}
\usepackage{parskip}
\usepackage{fancyhdr}
\usepackage{array}
\usepackage{booktabs}
\usepackage{enumitem}
\usepackage{amsmath}
\usepackage{tabularx}

\pagestyle{fancy}
\fancyhf{}
\cfoot{\thepage}
\rhead{Lecture 6}
\lhead{MA 214: Applied Statistics (Spring 2026)}
\renewcommand{\headrulewidth}{.4mm}

\newcommand{\blankline}[1]{\underline{\hspace{#1}}}

\begin{document}

\noindent\textbf{Name:}\blankline{2.25in}\hfill\textbf{BUID:}\blankline{1.55in}\newline

\begin{center}
 \Large In-Class Worksheet: Logistic Regression
\end{center}

\subsection*{Scenario}
Logistic regression is used when the response variable is binary. Instead of modeling the raw 0/1 outcome with a line, we model the probability of success using odds, log odds, and maximum likelihood.

\medskip
\noindent\textbf{Goals:} \textbf{(1)} identify binary responses and success coding, \textbf{(2)} convert between probabilities, odds, and log odds, \textbf{(3)} interpret logistic regression coefficients using odds ratios, \textbf{(4)} understand classification thresholds and unbalanced data.

\subsection*{Part 1: Binary Response and Success Coding}
For each question, define a binary response variable $Y$ and choose what value will count as success ($Y=1$).

\begin{center}
\renewcommand{\arraystretch}{1.45}
\begin{tabularx}{\textwidth}{p{5.2cm}p{3.2cm}p{3.2cm}X}
\toprule
\textbf{Question} & \textbf{$Y=0$} & \textbf{$Y=1$} & \textbf{What is $p=P(Y=1)$?} \\
\midrule
Did a passenger survive? & & & \\[1.8em]
Is an email spam? & & & \\[1.8em]
Did a borrower default? & & & \\
\bottomrule
\end{tabularx}
\end{center}

\begin{enumerate}[label={\bfseries (\Alph*)},itemsep=.8em]
\item Why is ordinary linear regression often a poor fit for a 0/1 response? \\[3.5em]
\end{enumerate}

\subsection*{Part 2: Probability, Odds, and Log Odds}
For an event with probability $p$,
\[
\text{odds}=\frac{p}{1-p}, \qquad \text{log odds}=\log\left(\frac{p}{1-p}\right).
\]
Complete the missing values.

\begin{center}
\renewcommand{\arraystretch}{1.6}
\begin{tabular}{ccc}
\toprule
\boldmath$p$ & \textbf{odds} & \textbf{log odds} \\
\midrule
0.50 & & \\[1.8em]
0.80 & & \\[1.8em]
0.25 & & \\
\bottomrule
\end{tabular}
\end{center}

\begin{enumerate}[label={\bfseries (\Alph*)},itemsep=.8em]
\item If the odds of survival are 3, is survival more likely or less likely than death? Explain briefly. \\[3em]
\end{enumerate}

\subsection*{Part 3: Interpreting a Logistic Regression Model}
For the Donner Party example, code $Y=1$ if a person survived. Using Male as the reference level, a fitted logistic model is
\[
\log\left(\frac{p}{1-p}\right) = 1.633 - 0.078(\text{Age}) + 1.597(\text{Female}),
\]
where $p$ is the probability of survival.

\begin{enumerate}[label={\bfseries (\Alph*)},itemsep=.8em]
\item What is the reference group for sex? \\[2em]
\item Convert the age coefficient into an odds ratio.
\[
e^{-0.078}=\blankline{2in}
\]
Interpret it in context. \\[3.5em]
\item Convert the female coefficient into an odds ratio.
\[
e^{1.597}=\blankline{2in}
\]
Interpret it in context. \\[3.5em]
\item For a 25-year-old male, compute the linear predictor $\eta$.
\[
\eta = 1.633 - 0.078(25) + 1.597(0)=\blankline{2in}
\]
\item Convert this $\eta$ to a probability using $p=e^\eta/(1+e^\eta)$.
\[
p = \blankline{2.5in}
\]
\end{enumerate}

\subsection*{Part 4: Likelihood and Classification}
For one binary observation, the Bernoulli likelihood contribution is
\[
P(Y_i=y_i)=p_i^{y_i}(1-p_i)^{1-y_i}.
\]

\begin{enumerate}[label={\bfseries (\Alph*)},itemsep=.8em]
\item If $y_i=1$, what is the likelihood contribution? \hfill \blankline{1.5in}
\item If $y_i=0$, what is the likelihood contribution? \hfill \blankline{1.5in}
\item A logistic model gives a probability $\hat p$, not automatically a class label. If we use threshold 0.50, fill in the rule:
\[
\hat p \ge 0.50 \Rightarrow \hat Y = \blankline{.7in}, \qquad \hat p < 0.50 \Rightarrow \hat Y = \blankline{.7in}.
\]
\end{enumerate}

\subsection*{Part 5: Misleading Accuracy}
Out of 1000 emails, 950 are not spam and 50 are spam. A classifier predicts every email is not spam.

\begin{center}
\renewcommand{\arraystretch}{1.35}
\begin{tabular}{lcc}
\toprule
 & \textbf{Actual not spam} & \textbf{Actual spam} \\
\midrule
Predicted not spam & 950 & 50 \\
Predicted spam & 0 & 0 \\
\bottomrule
\end{tabular}
\end{center}

\begin{enumerate}[label={\bfseries (\Alph*)},itemsep=.8em]
\item Compute the accuracy.
\[
\text{accuracy}=\blankline{2.5in}
\]
\item What is the problem with this classifier? \\[3.5em]
\item In one sentence, explain why a threshold other than 0.50 might be reasonable. \\[3.5em]
\end{enumerate}

\end{document}
